\section{Билет 17. Метод приведённого градиента Вульфа}

\ticketq{Метод приведённого градиента Вульфа.}

\paragraph{Постановка.}
\[
  \varphi(x) \to \min , \qquad Ax = b , \qquad x \ge 0 ,
\]
$\varphi$ выпукла и дифференцируема, $A$ размера $m\times n$, $m < n$,
$\rank A = m$.

\textbf{Условие невырожденности:} любые $m$ столбцов $A$ линейно независимы и
любая крайняя точка имеет ровно $m$ положительных компонент.

\textbf{Идея:} равенства используются для выражения базисных переменных через
небазисные; градиент пересчитывается с учётом этой зависимости; движение идёт
по антиприведённому градиенту.

\paragraph{Разбиение переменных.}
$x = (x_B, x_N)$, $A = (B,N)$, $B$ невырождена, $x_B > 0$:
\[
  Ax = b \iff Bx_B + Nx_N = b \iff x_B = B^{-1}(b - Nx_N) .
\]

\textbf{Лемма (сохранение равенств).} $Ax = b$, $Ad = 0$ $\Rightarrow$
$A(x+\alpha d) = Ax + \alpha Ad = b$. $\blacksquare$

Разбивая $d = (d_B, d_N)$, из $Ad = 0$ получаем $Bd_B + Nd_N = 0$, то есть
\[
  d_B = -B^{-1}N\,d_N :
\]
достаточно выбрать небазисную часть, базисная определяется однозначно.

\paragraph{Приведённый градиент.}
\[
  n\T = \nabla\T\varphi - \nabla_B\T\varphi\,B^{-1}A
      = \big(0,\ \nabla_N\T\varphi - \nabla_B\T\varphi\,B^{-1}N\big) ,
\]
то есть $n_N = \nabla_N\varphi - N\T B^{-\mathsf{T}}\nabla_B\varphi$;
компоненты обозначаем $n^{(j)}$, $j \in N$.

\textbf{Теорема.} \textit{Для допустимого направления ($Ad = 0$):}
$\nabla\T\varphi\,d = n_N\T d_N$.

\textsc{Док-во.} Подставим $d_B = -B^{-1}Nd_N$:
\[
  \nabla\T\varphi\,d
   = \nabla_B\T\varphi\,d_B + \nabla_N\T\varphi\,d_N
   = \big(\nabla_N\T\varphi - \nabla_B\T\varphi\,B^{-1}N\big)d_N
   = n_N\T d_N . \ \blacksquare
\]
Задача свелась к движению по небазисным переменным, и роль градиента в этом
пространстве играет $n_N$.

\paragraph{Направление Вульфа.}
Двигаться просто по $-n_N$ нельзя: при $x_j = 0$ отрицательная компонента
нарушила бы $x_j \ge 0$. Поэтому
\[
  d_j =
  \begin{cases}
    -\,n^{(j)} , & n^{(j)} \le 0 , \\[2pt]
    -\,x^{(j)}\,n^{(j)} , & n^{(j)} > 0 ,
  \end{cases}
  \qquad j \in N ,
  \qquad d_B = -B^{-1}N\,d_N .
\]
Множитель $x^{(j)}$ во втором случае и решает проблему: при $x_j = 0$
компонента $d_j$ обнуляется сама.

\textbf{Допустимость.} Для $j \in N$ с $x_j = 0$: либо $n^{(j)} \le 0$ и
$d_j = -n^{(j)} \ge 0$, либо $n^{(j)} > 0$ и $d_j = -x_jn^{(j)} = 0$ — в
обоих случаях $d_j \ge 0$. Базисные компоненты строго положительны
(невырожденность), поэтому при малом $\alpha$ имеем $x + \alpha d \ge 0$.

\paragraph{Теорема (направление спуска).}
\textit{$\nabla\T\varphi\,d \le 0$, причём равенство $\iff d_N = 0$.}

\textsc{Док-во.} По предыдущей теореме и правилу Вульфа:
\[
  n_N\T d_N
   = -\sum_{n^{(j)} \le 0}\big(n^{(j)}\big)^2
     - \sum_{n^{(j)} > 0} x^{(j)}\big(n^{(j)}\big)^2 \ \le\ 0 ,
\]
так как $x^{(j)} \ge 0$. Равенство требует обнуления всех слагаемых, то есть
\[
  n^{(j)} \ge 0 \ \ \forall j \in N
  \qquad\text{и}\qquad
  n^{(j)}x^{(j)} = 0 \ \ \forall j \in N . \ \blacksquare
\]
Это и есть \textbf{условия Куна--Таккера} ($z = n$: $z \ge 0$, $z\T x = 0$,
$z_B = 0$); при выпуклой $\varphi$ они достаточны для глобального минимума.
Кроме того, $d_N = 0 \Rightarrow d_B = 0$, то есть останов при $d = 0$
корректен.

\paragraph{Выбор шага.}
Равенства сохраняются автоматически; нужна неотрицательность:
\[
  \alpha_{\max} = \min_{i\,:\ d_i < 0}\left(-\frac{x_i}{d_i}\right)
  \qquad (+\infty, \text{ если все } d_i \ge 0) ,
\]
\[
  \alpha_k = \argmin_{0 \le \alpha \le \alpha_{\max}}
    \varphi(x^k + \alpha d^k) .
\]
При $d_i \ge 0$ очевидно $x_i + \alpha d_i \ge 0$; при $d_i < 0$ из
$\alpha \le x_i/(-d_i)$ следует $\alpha(-d_i) \le x_i$, то есть
$x_i + \alpha d_i \ge 0$.

\paragraph{Алгоритм.}
\begin{csteps}
  \item Задать допустимую $x^0 \ge 0$, $Ax^0 = b$, выбрать базис $B$.
  \item Вычислить $\nabla\varphi(x^k)$, приведённый градиент $n_N$, затем
        $d_N$ по правилу Вульфа и $d_B = -B^{-1}Nd_N$.
  \item $d = 0$ — конец (в выпуклом случае точка оптимальна).
  \item Найти $\alpha_{\max}$ и $\alpha_k$; положить
        $x^{k+1} = x^k + \alpha_k d$.
  \item Если базисная переменная обратилась в нуль — заменить её небазисной с
        наибольшей компонентой, пересчитать $B$. К~шагу~2.
\end{csteps}

\paragraph{Теоремы о сходимости.}
\textit{Последовательность $\varphi(x^k)$ убывает; любая предельная точка
$\{x^k\}$ — оптимальное решение.}

\textsc{Схема.} $\varphi(x^k)$ убывает и ограничена снизу $\Rightarrow$
сходится. Базисов конечное число, поэтому берём подпоследовательность с
фиксированным базисом; при фиксированном базисе направление $d(x)$ непрерывно
по $x$. Если для предельной точки $\bar{x}$ было бы $d(\bar{x}) \ne 0$, то
нашёлся бы шаг $\bar\alpha > 0$ с
$\varphi(\bar{x}+\bar\alpha d(\bar{x})) < \varphi(\bar{x})$, что противоречит
сходимости $\varphi(x^k)$. Значит $d(\bar{x}) = 0$: точка стационарна, а при
выпуклой $\varphi$ оптимальна. $\blacksquare$

\paragraph{Ограничения.}
Нужна допустимая начальная точка; требуется выбирать и обновлять базис; при
вырождении возможны трудности; плохо обусловленная $B$ даёт неустойчивость;
при $\alpha_k = \alpha_{\max}$ базис меняется; для невыпуклой $\varphi$ —
только локальный минимум.

\seefull{гл.~3, \S14 (численный пример там же)}
